\section{Supplementary Material}\label{supplemental}

Figure~\ref{fig:qmc-dft-cdens} shows FP-DMC, HF, and DFT charge densities before and after the metal-insulator transition (MIT) at $r_s=3$. Hybrid LDA is the only density functional that reasonably reproduces QMC charge density on both sides of the MIT. In Fig.~\ref{fig:dft-mit}, we show the evolution of charge density across the MIT. Only hybrid LDA and QMC have a discontinuous jump across the transition.

\begin{figure}[h]
\begin{minipage}{0.48\textwidth}
\includegraphics[width=\linewidth]{tb35_nc-para-rs3.00-mu0.300-lda-n144-cdens.pdf}
(a) metal $V_M/W=0.3$
\end{minipage}
\begin{minipage}{0.48\textwidth}
\includegraphics[width=\linewidth]{tb35_nc-120-rs3.00-mu0.400-hf-n144-cdens.pdf}
(b) insulator $V_M/W=0.4$
\end{minipage}
\caption{FP-DMC (N=144) charge density at $r_s=3$ along the linecut BAC. Charge densities from three effective one-body theories are shown along with the QMC linecuts. All functionals produce reasonable charge density for the metallic phase, but only hybrid LDA agrees well with QMC in the insulating phase.}
\label{fig:qmc-dft-cdens}
\end{figure}

\begin{figure}[h]
\begin{minipage}{0.48\textwidth}
\includegraphics[width=\linewidth]{tb35_nc-120-rs3.00-hf-dft-cdens}
(a) HF
\end{minipage}
\begin{minipage}{0.48\textwidth}
\includegraphics[width=\linewidth]{tb35_nc-120-rs3.00-lda-dft-cdens}
(b) LDA
\end{minipage}
\begin{minipage}{0.48\textwidth}
\includegraphics[width=\linewidth]{tb35_nc-120-rs3.00-exx0.50-dft-cdens}
(c) LDA+$50$\%EXX
\end{minipage}
\begin{minipage}{0.48\textwidth}
\includegraphics[width=\linewidth]{tb35_nc-120-rs3.00-qmc-dft-cdens}
(d) QMC
\end{minipage}
\caption{Charge density linecuts (same as Fig.~\ref{fig:qmc-dft-cdens}) across the metal insulator transition at $r_s=3$.}
\label{fig:dft-mit}
\end{figure}

Table~\ref{tab:dmc-energy} lists the FP-DMC energy for the calculations shown in Fig.~\ref{fig:qmc-dft-cdens}. At each point on the phase diagram, we compute two FP-DMC energies using Slater-Jastrow trial wavefunctions built from HF and LDA orbitals. We choose the lowest-energy trial to calculate QMC properties. The LDA trial changes from a magnetic insulator to a magnetic metal (spin density wave phase) then to a non-magnetic metal. The HF trial is magnetic and insulating and its FP-DMC energy becomes higher than that of the LDA trial in only the non-magnetic metal phase.

\begin{table}[h]
\caption{Twist averaged FP-DMC energies in Hartree at $r_s=3$.}
\label{tab:dmc-energy}
\begin{minipage}{0.48\textwidth}
(a) metallic phase \\ % tb35/nc
\begin{tabular}{lrrrll}
\toprule
   mag &  $r_s$ &  $V_M/W$ &   N & func &         E/N \\
\midrule
   120 &    3.0 &      0.3 & 144 &   HF & -0.21824(6) \\
   120 &    3.0 &      0.3 & 144 &  LDA & -0.22130(1) \\
  para &    3.0 &      0.3 & 144 &   HF & -0.22121(1) \\
  para &    3.0 &      0.3 & 144 &  LDA & -0.22129(1) \\
stripe &    3.0 &      0.3 & 144 &   HF & -0.21828(5) \\
stripe &    3.0 &      0.3 & 144 &  LDA & -0.22128(1) \\
\bottomrule
\end{tabular}
\end{minipage}
\begin{minipage}{0.48\textwidth}
(b) insulating phase \\
\begin{tabular}{lrrrll}
\toprule
   mag &  $r_s$ &  $V_M/W$ &   N & func &         E/N \\
\midrule
   120 &    3.0 &      0.4 & 144 &   HF & -0.23262(4) \\
   120 &    3.0 &      0.4 & 144 &  LDA & -0.23003(2) \\
  para &    3.0 &      0.4 & 144 &   HF & -0.23022(2) \\
  para &    3.0 &      0.4 & 144 &  LDA & -0.23037(2) \\
stripe &    3.0 &      0.4 & 144 &   HF & -0.23240(4) \\
stripe &    3.0 &      0.4 & 144 &  LDA & -0.23039(2) \\
\bottomrule
\end{tabular}
\end{minipage}
\end{table}

Figure~\ref{fig:daf-dens} shows a comparison of FP-DMC and ph-AFQMC magnetization densities of the $120^\circ$ N\'eel phase at $r_s=3$ and $\lambda=0.6$.
Both calculations averaged over $16$ twists, each containing $36$ electrons. On average, $140$ HF orbitals at each twist were used as basis functions in AFQMC.
The DMC calculation of the same system used a Slater-Jastrow trial wavefunction built from the occupied HF orbitals.
Densities from the two methods are consistent despite the very different approximations in the algorithms.

\begin{figure}[h]
% tbeg/035-ni-exx/nc/analysis/reply_daf_dens.py
\begin{minipage}{0.48\textwidth}
\includegraphics[width=\linewidth]{{tb35_nc-120-rs3.00-mu0.600-n36-x-sdens2d}}
(a) $\rho_x(x, y)$
\end{minipage}
\begin{minipage}{0.48\textwidth}
\includegraphics[width=\linewidth]{{tb35_nc-120-rs3.00-mu0.600-n36-y-sdens2d}}
(b) $\rho_y(x, y)$
\end{minipage}
\begin{minipage}{0.48\textwidth}
\includegraphics[width=\linewidth]{{tb35_nc-120-rs3.00-mu0.600-n36-mx-sdens1d}}
(c) $\rho_x(r, r)$
\end{minipage}
\begin{minipage}{0.48\textwidth}
\includegraphics[width=\linewidth]{{tb35_nc-120-rs3.00-mu0.600-n36-my-sdens1d}}
(d) $\rho_y(r, r)$
\end{minipage}
\caption{Magnetization density of AFQMC and DMC at $r_s=3$ and $\lambda=0.6$.
}
\label{fig:daf-dens}
\end{figure}

To make an order-of-magnitute estimate for the magnetic ordering temperature, we calculate the energy cost of a single spin flip $\Delta E$ and related it to the exchange parameter $J=\Delta E/4$ in a nearest-neighbor Heisenberg model of the spin stripe phase.
We made the estimates at $r_s=8$ using a $36$-electron system and FP-DMC with Slater-Jastrow wavefunction and LDA orbitals.
%We made no attempt to correct systematic errors in these order-of-magnitude estimates.
As shown in Fig.~\ref{fig:jheis}, $J$ decreases rapidly from $O(100)$ mK at $\lambda=0.4$ to within errorbars of zero at $\lambda=1.6$.
Given the deep moir\'e potential measured in experiments~\cite{Shabani2021}, we do not expect the ground-state magnetic order to be measurable because of large thermal fluctuations.

\begin{figure}[h]
% tbeg/034-j12-heis/qmcco/gather/jheis.py
\begin{minipage}{0.48\textwidth}
\includegraphics[width=\linewidth]{{tb34_qmcco-gather-jheis}}
\end{minipage}
\caption{Estimates of the exchange parameter $J$ of a nearest-neighbor Heisenberg model at $r_s=8$.
}
\label{fig:jheis}
\end{figure}

Figure~\ref{fig:nc-extrap} shows finite-size correction (FSC)~\cite{Chiesa2006,Holzmann2016-fsc} and thermodynamic extrapolation~\cite{Drummond2009} at $r_s=3$ and $\lambda=0.3, 0.4$.
The two approaches agree in simulation cells containing $144$ electrons in both the metallic and the insulating phases.
To calculate the FSC, we fit the fluctuating structure factor $\delta S(k)=\langle (\rho_{\bs{k}}-\langle\rho_{\bs{k}}\rangle)(\rho_{-\bs{k}}-\langle\rho_{-\bs{k}}\rangle)\rangle$ to $a k^{3/2} + b k^2 + c k^4$ as shown in Fig.~\ref{fig:nc-dsk}.
The potential energy correction is calculated as $\delta V= [\int -\sum] \frac{1}{2}v_k \delta S(k)$, while the kinetic energy correction is approximated to leading-order $\delta T \approx \delta V$.
These results show that the $144$-electron system is large enough to accurately estimate the total energy in the thermodynamic limit.

\begin{figure}[h]
% tbeg/035-ni-exx/nc/analysis/nc_extrap.py
\begin{minipage}{0.48\textwidth}
\includegraphics[width=\linewidth]{{tb35_nc-extrap-rs3-mu0.3}}
(a) $V_M/W=0.3$
\end{minipage}
\begin{minipage}{0.48\textwidth}
\includegraphics[width=\linewidth]{{tb35_nc-extrap-rs3-mu0.4}}
(b) $V_M/W=0.4$
\end{minipage}
\caption{Finite-size correction (FSC) and extrapolation of FP-DMC total energy at $r_s=3$.
The faint dashed line shows size extrapolation.
The horizontal line and its width represent the mean and error of the thermodynamic limit.
It agrees with $N=144$ result with FSC.
}
\label{fig:nc-extrap}
\end{figure}

\begin{figure}[h]
% tbeg/035-ni-exx/nc/analysis/fsc_etot.py
\begin{minipage}{0.48\textwidth}
\includegraphics[width=\linewidth]{{tb35_nc-dsk-fsc-para-rs3-mu0.3-lda}}
(a) non-magnetic LDA orbitals $V_M/W=0.3$
\end{minipage}
\begin{minipage}{0.48\textwidth}
\includegraphics[width=\linewidth]{{tb35_nc-dsk-fsc-120-rs3-mu0.3-hf}}
(b) $120^\circ$ N\'eel HF orbitals $V_M/W=0.3$
\end{minipage}
\begin{minipage}{0.48\textwidth}
\includegraphics[width=\linewidth]{{tb35_nc-dsk-fsc-para-rs3-mu0.4-lda}}
(c) non-magnetic LDA orbitals $V_M/W=0.4$
\end{minipage}
\begin{minipage}{0.48\textwidth}
\includegraphics[width=\linewidth]{{tb35_nc-dsk-fsc-120-rs3-mu0.4-hf}}
(d) $120^\circ$ N\'eel HF orbitals $V_M/W=0.4$
\end{minipage}
\caption{Fluctuating structure factor $\delta S(k)$ at $r_s=3$.
The black solid line is reference from random phase approximation (RPA), which exactly captures the long-wavelength plasmon fluctuations in the metallic phase.
Blue lines are polynomial fits to the FP-DMC data.
}
\label{fig:nc-dsk}
\end{figure}
